% Generated draft: requires mathematical and source review.
\documentclass[11pt]{article}
\usepackage[margin=25mm]{geometry}
\usepackage{fontspec}
\setmainfont{Latin Modern Roman}
\usepackage{amsmath,amssymb,mathtools,unicode-math}
\setmathfont{Latin Modern Math}
\usepackage{enumitem,needspace,xurl,hyperref}
\setlength{\parindent}{0pt}
\setlength{\parskip}{5pt}
\setlength{\emergencystretch}{3em}
\newcommand{\N}{\mathbb{N}}
\newcommand{\Z}{\mathbb{Z}}
\newcommand{\Q}{\mathbb{Q}}
\newcommand{\R}{\mathbb{R}}
\newcommand{\C}{\mathbb{C}}
\newcommand{\F}{\mathbb{F}}
\newcommand{\A}{\mathbb{A}}
\newcommand{\PP}{\mathbb{P}}
\newcommand{\E}{\mathbb{E}}
\newcommand{\eps}{\varepsilon}
\newcommand{\dd}{\,\mathrm{d}}
\begin{document}
\section*{UnsolvedMath ID: 30002141}
\textbf{Source problem code:} catalogue entry
\subsection{Definitions and mathematical statement}
The supplied question concerns compact negatively curved surfaces \(M_1,M_2\) and a homotopy equivalence \(f:M_1\to M_2\). It does not explicitly state Thurston's length-spectrum criterion, the permitted values of \(L\), or the boundary hypotheses. Consequently, the following is a conditional reconstruction of the intended question, requiring confirmation against the source.

Under the Riemannian interpretation, let \(g_i\) be the negatively curved metric on \(M_i\), let \(d_i\) be its intrinsic distance, and, for a free homotopy class \(\alpha\) of loops in \(M_i\), define
\[
\ell_i(\alpha)=\inf\{\operatorname{Length}_{g_i}(\gamma):\gamma\text{ is a piecewise smooth loop representing }\alpha\}.
\]
The map \(f\) induces a map \(f_*\) on free homotopy classes by composition. Being a homotopy equivalence means that there is a continuous map \(q:M_2\to M_1\) with \(q\circ f\simeq\operatorname{id}_{M_1}\) and \(f\circ q\simeq\operatorname{id}_{M_2}\).

For a prescribed constant \(L>0\), a natural interpretation of the unspecified length-spectrum hypothesis is
\[
\ell_2(f_*\alpha)\leq L\,\ell_1(\alpha)
\qquad\text{for every nontrivial free homotopy class }\alpha\text{ in }M_1.
\tag{LS}
\]
The corresponding realization question is whether, for every pair of surfaces satisfying the intended hypotheses and every homotopy equivalence \(f\) satisfying \((\mathrm{LS})\), there exists a homeomorphism \(h:M_1\to M_2\) such that
\[
h\simeq f,
\qquad
d_2(h(x),h(y))\leq L\,d_1(x,y)
\quad\text{for all }x,y\in M_1.
\]
The constant in the conclusion is the same prescribed \(L\) as in the length comparison.

The supplied record does not establish that \((\mathrm{LS})\), with all nontrivial free homotopy classes and all \(L>0\), is the exact intended criterion. In particular, a restriction to simple closed curves, a restriction such as \(L\geq1\), or additional geometric hypotheses must not be silently assumed. It is also unspecified whether the surfaces are closed or may have boundary; in the latter case, boundary geometry and any requirements on maps, homotopies, and boundary lengths need clarification. Connectedness, orientability, metric regularity, and the precise curvature assumptions beyond the phrase ``negatively curved'' are not explicitly stated.
\subsection{Short English statement}
The intended question asks whether a marked comparison of closed-curve lengths guarantees a homeomorphism in the given homotopy class whose Lipschitz constant is at most the same prescribed \(L\). The displayed formulation is provisional because the record does not spell out the criterion, the range of \(L\), or the boundary hypotheses. The supplied status is partially solved; this draft does not independently verify that assessment.
\subsection{Sources}
\begin{itemize}
\item Supplied local UnsolvedMath registry record: permanent ID \(30002141\), exact source problem code \texttt{catalogue entry}, title ``Lipschitz Realization of Surface Homotopy Equivalences,'' in the Oberwolfach Reports Open Problems collection. The record gives proposed year \(2012\) and names no proposer.
\item Supplied report citation: \emph{Geometric Group Theory, Hyperbolic Dynamics and Symplectic Geometry} (2013), Oberwolfach Reports, pp.~2139--2203, DOI \texttt{10.4171/owr/2012/35}. \url{https://doi.org/10.4171/owr/2012/35}. This is also the supplied literature link labelled ``Oberwolfach Report 35.'' The citation year is retained as supplied, alongside the record's proposed year \(2012\).
\item Supplied literature reference: Butt, ``quantitative marked length spectrum rigidity.'' \url{https://arxiv.org/abs/2203.12128}.
\item Supplied literature reference: Thurston, ``minimal stretch maps.'' \url{https://arxiv.org/abs/math/9801039}.
\item Literature assessment dated exactly \texttt{2026-08-21}: the supplied curation labels the problem partially solved and reports quantitative local marked-length-spectrum rigidity for negatively curved surfaces and related locally symmetric cases, while describing the full extension of Thurston's global theorem as unresolved. This is supplied dated curation, not an independently verified literature assessment; the background expressly requests independent verification of time-sensitive status and cited claims.
\item The registry labels the statement \texttt{corrected\_verified} and reports that its mathematical content was checked against the official Oberwolfach report, with notation, context, and question form normalized without changing the source-level claim. That verification assertion is supplied provenance and has not been independently checked in this draft.
\end{itemize}
\end{document}
